\documentclass[10pt,a4paper]{amsart}
\usepackage{amsmath,amssymb,amsthm,graphicx}
\usepackage{verbatim}
\newenvironment{reference}{\comment}{\endcomment}
%\newenvironment{reference}{}{}
\newenvironment{slogan}{\comment}{\endcomment}
\newenvironment{history}{\comment}{\endcomment}

% For commutative diagrams we use Xy-pic
\usepackage[all]{xy}

% We use 2cell for 2-commutative diagrams.
\xyoption{2cell}
% Theorem environments.
%
\theoremstyle{plain}
\newtheorem{theorem}[subsection]{Theorem}
\newtheorem{proposition}[subsection]{Proposition}
\newtheorem{lemma}[subsection]{Lemma}

\theoremstyle{definition}
\newtheorem{definition}[subsection]{Definition}
\newtheorem{example}[subsection]{Example}
\newtheorem{exercise}[subsection]{Exercise}
\newtheorem{situation}[subsection]{Situation}

\theoremstyle{remark}
\newtheorem{remark}[subsection]{Remark}
\newtheorem{remarks}[subsection]{Remarks}

\numberwithin{equation}{subsection}

% Macros
%
\def\lim{\mathop{\mathrm{lim}}\nolimits}
\def\colim{\mathop{\mathrm{colim}}\nolimits}
\def\Spec{\mathop{\mathrm{Spec}}}
\def\Hom{\mathop{\mathrm{Hom}}\nolimits}
\def\Ext{\mathop{\mathrm{Ext}}\nolimits}
\def\SheafHom{\mathop{\mathcal{H}\!\mathit{om}}\nolimits}
\def\SheafExt{\mathop{\mathcal{E}\!\mathit{xt}}\nolimits}
\def\Sch{\mathit{Sch}}
\def\Mor{\mathop{\mathrm{Mor}}\nolimits}
\def\Ob{\mathop{\mathrm{Ob}}\nolimits}
\def\Sh{\mathop{\mathit{Sh}}\nolimits}
\def\NL{\mathop{N\!L}\nolimits}
\def\CH{\mathop{\mathrm{CH}}\nolimits}
\def\proetale{{pro\text{-}\acute{e}tale}}
\def\etale{{\acute{e}tale}}
\def\QCoh{\mathit{QCoh}}
\def\Ker{\mathop{\mathrm{Ker}}}
\def\Im{\mathop{\mathrm{Im}}}
\def\Coker{\mathop{\mathrm{Coker}}}
\def\Coim{\mathop{\mathrm{Coim}}}

% Boxtimes
%
\DeclareMathSymbol{\boxtimes}{\mathbin}{AMSa}{"02}

%
% Macros for moduli stacks/spaces
%
\def\QCohstack{\mathcal{QC}\!\mathit{oh}}
\def\Cohstack{\mathcal{C}\!\mathit{oh}}
\def\Spacesstack{\mathcal{S}\!\mathit{paces}}
\def\Quotfunctor{\mathrm{Quot}}
\def\Hilbfunctor{\mathrm{Hilb}}
\def\Curvesstack{\mathcal{C}\!\mathit{urves}}
\def\Polarizedstack{\mathcal{P}\!\mathit{olarized}}
\def\Complexesstack{\mathcal{C}\!\mathit{omplexes}}
% \Pic is the operator that assigns to X its picard group, usage \Pic(X)
% \Picardstack_{X/B} denotes the Picard stack of X over B
% \Picardfunctor_{X/B} denotes the Picard functor of X over B
\def\Pic{\mathop{\mathrm{Pic}}\nolimits}
\def\Picardstack{\mathcal{P}\!\mathit{ic}}
\def\Picardfunctor{\mathrm{Pic}}
\def\Deformationcategory{\mathcal{D}\!\mathit{ef}}
\usepackage{cleveref}
\crefname{originallemma}{original lemma}{original lemmas}
\Crefname{originallemma}{Original lemma}{Original lemmas}
\makeatletter
\newlabel{decent-spaces-lemma-decent-generic-points}{{0ABV}{source}{}{sourcepointer-1}{}}
\newlabel{decent-spaces-lemma-decent-generic-points@cref}{{[originallemma][0][]0ABV}{[1][0][]source}{}{}{}}
\makeatother
\crefname{originallemma}{original lemma}{original lemmas}
\Crefname{originallemma}{Original lemma}{Original lemmas}
\makeatletter
\newlabel{decent-spaces-lemma-locally-Noetherian-decent-quasi-separated}{{0BB6}{source}{}{sourcepointer-2}{}}
\newlabel{decent-spaces-lemma-locally-Noetherian-decent-quasi-separated@cref}{{[originallemma][0][]0BB6}{[1][0][]source}{}{}{}}
\makeatother
\crefname{originalreference}{original reference}{original references}
\Crefname{originalreference}{Original reference}{Original references}
\makeatletter
\newlabel{definition-integral-algebraic-space}{{0AD4}{source}{}{sourcepointer-3}{}}
\newlabel{definition-integral-algebraic-space@cref}{{[originalreference][0][]0AD4}{[1][0][]source}{}{}{}}
\makeatother
\crefname{originalreference}{original reference}{original references}
\Crefname{originalreference}{Original reference}{Original references}
\makeatletter
\newlabel{lemma-alteration-contained-in}{{0DMN}{source}{}{sourcepointer-4}{}}
\newlabel{lemma-alteration-contained-in@cref}{{[originalreference][0][]0DMN}{[1][0][]source}{}{}{}}
\makeatother
\crefname{originalreference}{original reference}{original references}
\Crefname{originalreference}{Original reference}{Original references}
\makeatletter
\newlabel{lemma-integral-dimension}{{0ECG}{source}{}{sourcepointer-5}{}}
\newlabel{lemma-integral-dimension@cref}{{[originalreference][0][]0ECG}{[1][0][]source}{}{}{}}
\makeatother
\crefname{originallemma}{original lemma}{original lemmas}
\Crefname{originallemma}{Original lemma}{Original lemmas}
\makeatletter
\newlabel{morphisms-lemma-alteration-dimension}{{02JX}{source}{}{sourcepointer-6}{}}
\newlabel{morphisms-lemma-alteration-dimension@cref}{{[originallemma][0][]02JX}{[1][0][]source}{}{}{}}
\makeatother
\crefname{originalproposition}{original proposition}{original propositions}
\Crefname{originalproposition}{Original proposition}{Original propositions}
\makeatletter
\newlabel{spaces-limits-proposition-there-is-a-scheme-finite-over}{{09YC}{source}{}{sourcepointer-7}{}}
\newlabel{spaces-limits-proposition-there-is-a-scheme-finite-over@cref}{{[originalproposition][0][]09YC}{[1][0][]source}{}{}{}}
\makeatother
\crefname{originallemma}{original lemma}{original lemmas}
\Crefname{originallemma}{Original lemma}{Original lemmas}
\makeatletter
\newlabel{spaces-morphisms-lemma-dimension-formula-general}{{0BAX}{source}{}{sourcepointer-8}{}}
\newlabel{spaces-morphisms-lemma-dimension-formula-general@cref}{{[originallemma][0][]0BAX}{[1][0][]source}{}{}{}}
\makeatother
\crefname{originallemma}{original lemma}{original lemmas}
\Crefname{originallemma}{Original lemma}{Original lemmas}
\makeatletter
\newlabel{spaces-morphisms-lemma-locally-finite-type-locally-noetherian}{{04ZK}{source}{}{sourcepointer-9}{}}
\newlabel{spaces-morphisms-lemma-locally-finite-type-locally-noetherian@cref}{{[originallemma][0][]04ZK}{[1][0][]source}{}{}{}}
\makeatother
\crefname{originallemma}{original lemma}{original lemmas}
\Crefname{originallemma}{Original lemma}{Original lemmas}
\makeatletter
\newlabel{spaces-properties-lemma-dimension}{{0BAN}{source}{}{sourcepointer-10}{}}
\newlabel{spaces-properties-lemma-dimension@cref}{{[originallemma][0][]0BAN}{[1][0][]source}{}{}{}}
\makeatother
\title{The dimension formula for a dominant locally finite type morphism of integral algebraic spaces}
\author{The Stacks Project Authors}
\date{stacks-project:spaces-over-fields.tex:local-snapshot-20261012; GNU FDL 1.2 or later}
\begin{document}
\maketitle
\noindent\textit{Editorial note.} All delivered mathematical text is the author's own, copied from the pinned Stacks Project chapter file spaces-over-fields.tex; no statement or proof step was rewritten, reordered or removed.\par\medskip
\noindent\textit{Editorial note.} The counted proof is the author's complete argument: the derivation of the inequality from the pointwise dimension formula and the description of dimension as a supremum of local ring dimensions, the reduction to a quasi-compact situation with the production of finite surjective integral schemes over both spaces, and the comparison of dimensions and transcendence degrees reducing the equality to the case of schemes.\par\medskip
\noindent\textit{Editorial note.} The label lemma-alteration-dimension names the counted statement, and the proof environment immediately following it in the source is delivered verbatim as its proof.\par\medskip
\noindent\textit{Editorial note.} The counted statement is carried by the author's own lemma environment in the chapter source.\par\medskip
\noindent\textit{Editorial note.} Pre-open gate: python operations/agent-authoring/tools/precard\_gate.py --source spaces-over-fields.tex --label lemma-alteration-dimension --start 2810 --end 2887 --local-id 28056 returned RESULT=OK (L28056 free, occupancy FREE, public interval clear). Candidate taken from sources/stacks-spaces-over-fields/candidates-halfA.json.\par\medskip
\noindent\textit{Editorial note.} The statement and proof are copied from the pinned chapter file; the cited inputs are kept as external pointers with their canonical Stacks tags.\par\medskip
\noindent\textit{Allowed standard prerequisites (editorial).}\begin{itemize}
\item Dimension formula for morphisms of algebraic spaces. Morphisms of Spaces, the lemma bounding the dimension of the local ring at a point by the dimension at its image plus the transcendence degree of the residue field extension. Supplies the pointwise inequality used to deduce the inequality on dimensions.
\item Dimension as the supremum of the dimensions of the local rings. Properties of Spaces, the lemma computing the dimension of an algebraic space as the supremum of the dimensions of its local rings, and Same chapter, the description of the generic point of an integral algebraic space. Supply the identification of the dimensions as suprema and the fact that the generic point is the unique point with zero dimensional local ring.
\item Existence of a finite surjective scheme over a locally Noetherian decent space. Limits of Spaces, the proposition producing a scheme finite and surjective over a given algebraic space, together with Same chapter, the lemma producing an integral closed subspace contained in a given one. Supplies the integral schemes Y' and X'' finite and surjective over Y and over a closed subspace of the fibre product.
\item Invariance of dimension under a finite surjective morphism and the scheme case of the formula. Same chapter, the lemma that the dimension is unchanged under an integral surjective morphism, and Morphisms, the dimension formula for the case of schemes. Supply the equality of dimensions of X'' with X and of Y' with Y, and the conclusion from the case of schemes.
\end{itemize}\medskip
\ifdefined\hypertarget\hypertarget{sourcepointer-1}{}\fi
\noindent\textit{Source reference.} Original lemma 0ABV: sources/stacks-chow/tags.txt, source lines 13033--13033. Prerequisite pointer invoked at proof line(s) 2833; it names a standard fact of a sibling chapter that the counted argument uses, and every step of the counted proof built on it is retained verbatim in the delivered body.\par\medskip
\ifdefined\hypertarget\hypertarget{sourcepointer-2}{}\fi
\noindent\textit{Source reference.} Original lemma 0BB6: sources/stacks-chow/tags.txt, source lines 14234--14234. Prerequisite pointer invoked at proof line(s) 2853; it names a standard fact of a sibling chapter that the counted argument uses, and every step of the counted proof built on it is retained verbatim in the delivered body.\par\medskip
\ifdefined\hypertarget\hypertarget{sourcepointer-3}{}\fi
\noindent\textit{Source reference.} Original reference 0AD4: sources/stacks-spaces-over-fields/pinned-original-source.tex, source lines 438--438. Prerequisite pointer invoked at proof line(s) 2829; it names a standard fact of the same chapter that the counted argument uses, and every step of the counted proof built on it is retained verbatim in the delivered body.\par\medskip
\ifdefined\hypertarget\hypertarget{sourcepointer-4}{}\fi
\noindent\textit{Source reference.} Original reference 0DMN: sources/stacks-spaces-over-fields/pinned-original-source.tex, source lines 1414--1414. Prerequisite pointer invoked at proof line(s) 2859; it names a standard fact of the same chapter that the counted argument uses, and every step of the counted proof built on it is retained verbatim in the delivered body.\par\medskip
\ifdefined\hypertarget\hypertarget{sourcepointer-5}{}\fi
\noindent\textit{Source reference.} Original reference 0ECG: sources/stacks-spaces-over-fields/pinned-original-source.tex, source lines 2793--2793. Prerequisite pointer invoked at proof line(s) 2877; it names a standard fact of the same chapter that the counted argument uses, and every step of the counted proof built on it is retained verbatim in the delivered body.\par\medskip
\ifdefined\hypertarget\hypertarget{sourcepointer-6}{}\fi
\noindent\textit{Source reference.} Original lemma 02JX: sources/stacks-chow/tags.txt, source lines 3352--3352. Prerequisite pointer invoked at proof line(s) 2886; it names a standard fact of a sibling chapter that the counted argument uses, and every step of the counted proof built on it is retained verbatim in the delivered body.\par\medskip
\ifdefined\hypertarget\hypertarget{sourcepointer-7}{}\fi
\noindent\textit{Source reference.} Original proposition 09YC: sources/stacks-chow/tags.txt, source lines 12560--12560. Prerequisite pointer invoked at proof line(s) 2856; it names a standard fact of a sibling chapter that the counted argument uses, and every step of the counted proof built on it is retained verbatim in the delivered body.\par\medskip
\ifdefined\hypertarget\hypertarget{sourcepointer-8}{}\fi
\noindent\textit{Source reference.} Original lemma 0BAX: sources/stacks-chow/tags.txt, source lines 14225--14225. Prerequisite pointer invoked at proof line(s) 2835; it names a standard fact of a sibling chapter that the counted argument uses, and every step of the counted proof built on it is retained verbatim in the delivered body.\par\medskip
\ifdefined\hypertarget\hypertarget{sourcepointer-9}{}\fi
\noindent\textit{Source reference.} Original lemma 04ZK: sources/stacks-chow/tags.txt, source lines 6468--6468. Prerequisite pointer invoked at proof line(s) 2865; it names a standard fact of a sibling chapter that the counted argument uses, and every step of the counted proof built on it is retained verbatim in the delivered body.\par\medskip
\ifdefined\hypertarget\hypertarget{sourcepointer-10}{}\fi
\noindent\textit{Source reference.} Original lemma 0BAN: sources/stacks-chow/tags.txt, source lines 14216--14216. Prerequisite pointer invoked at proof line(s) 2841; it names a standard fact of a sibling chapter that the counted argument uses, and every step of the counted proof built on it is retained verbatim in the delivered body.\par\medskip
\begin{lemma}
\label{lemma-alteration-dimension}
Let $S$ be a scheme. Let $f : X \to Y$ be a morphism of algebraic spaces
over $S$. Assume that
\begin{enumerate}
\item $Y$ is locally Noetherian,
\item $X$ and $Y$ are integral algebraic spaces,
\item $f$ is dominant, and
\item $f$ is locally of finite type.
\end{enumerate}
If $x \in |X|$ and $y \in |Y|$ are the generic points, then
$$
\dim(X) \leq \dim(Y) + \text{transcendence degree of }x/y.
$$
If $f$ is proper, then equality holds.
\end{lemma}

\begin{proof}
Recall that $|X|$ and $|Y|$ are irreducible sober topological spaces, see
discussion following Definition \ref{definition-integral-algebraic-space}.
Thus the fact that $f$ is dominant means that $|f|$ maps $x$ to $y$.
Moreover, $x \in |X|$ is the unique point at which the
local ring of $X$ has dimension $0$, see
Decent Spaces, Lemma \ref{decent-spaces-lemma-decent-generic-points}.
By Morphisms of Spaces, Lemma
\ref{spaces-morphisms-lemma-dimension-formula-general}
we see that the dimension of the local ring of $X$ at
any point $x' \in |X|$ is at most the dimension of the local
ring of $Y$ at $y' = f(x')$ plus the transcendence degree of $x/y$.
Since the dimension of $X$, resp.\ dimension of $Y$ is the
supremum of the dimensions of the local rings at $x'$, resp.\ $y'$
(Properties of Spaces, Lemma \ref{spaces-properties-lemma-dimension})
we conclude the inequality holds.

\medskip\noindent
Assume $f$ is proper.
Let $V \subset Y$ be a nonempty quasi-compact open subspace.
If we can prove the equality for the morphism $f^{-1}(V) \to V$,
then we get the equality for $X \to Y$. Thus we may assume that
$X$ and $Y$ are quasi-compact.
Observe that $X$ is quasi-separated as
a locally Noetherian decent algebraic space, see
Decent Spaces, Lemma
\ref{decent-spaces-lemma-locally-Noetherian-decent-quasi-separated}.
Thus we may choose $Y' \to Y$ finite surjective where $Y'$
is a scheme, see Limits of Spaces, Proposition
\ref{spaces-limits-proposition-there-is-a-scheme-finite-over}.
After replacing $Y'$ by a suitable closed subscheme, we
may assume $Y'$ is integral, see for example the more general
Lemma \ref{lemma-alteration-contained-in}.
By the same lemma, we may choose a closed subspace
$X' \subset X \times_Y Y'$ such that $X'$ is integral
and $X' \to X$ is finite surjective.
Now $X'$ is also locally Noetherian
(Morphisms of Spaces, Lemma
\ref{spaces-morphisms-lemma-locally-finite-type-locally-noetherian})
and we can use Limits of Spaces, Proposition
\ref{spaces-limits-proposition-there-is-a-scheme-finite-over}
once more to choose a finite surjective morphism $X'' \to X'$
with $X''$ a scheme. As before we may assume that $X''$
is integral. Picture
$$
\xymatrix{
X'' \ar[d] \ar[r] & X \ar[d]^f \\
Y' \ar[r] & Y
}
$$
By Lemma \ref{lemma-integral-dimension} we have $\dim(X'') = \dim(X)$
and $\dim(Y') = \dim(Y)$. Since $X$ and $Y$ have open neighbourhoods
of $x$, resp.\ $y$ which are schemes, we readily see that the generic
points $x'' \in X''$, resp.\ $y' \in Y'$ are the unique points mapping
to $x$, resp.\ $y$ and that the residue field extensions
$\kappa(x'')/\kappa(x)$ and $\kappa(y')/\kappa(y)$ are finite.
This implies that the transcendence degree of $x''/y'$ is the
same as the transcendence degree of $x/y$. Thus the equality follows
from the case of schemes whicn is
Morphisms, Lemma \ref{morphisms-lemma-alteration-dimension}.
\end{proof}

\end{document}
